% ProveIt research insertion fragment. No document preamble.
% Requires the parent's standard theorem/proof environments and AMS math.
% Full proofs, scope, literature and checks: accompanying article.tex.
% All labels in this fragment are prefixed connected:.
\section{Connected spectral jets and mixed Gamma primitives}
\label{connected:section}
Let $a>0$, $x_n=n+a$, $m\ge1$, and $\operatorname{Re}(a+u_j)>0$.
Each linear factor uses its right-half-plane principal logarithm. A spectral
product uses the sum of those factor logarithms, not an assumed principal-log
multiplicativity. Put $U=\prod_j u_j$ and $S=\sum_j s_j$, and define
\[
 D_s(x,u)=-\int_0^u(x+t)^{-1-s}\,dt
 =\frac{(x+u)^{-s}-x^{-s}}s,
\]
where the quotient is continued at $s=0$. Thus $D_0=-\log(1+u/x)$.
The connected critical spectrum is
\[
 \mathcal C_m(\boldsymbol{s})
 =\sum_{n\ge0}x_n^{m-1}
    \prod_{j=1}^m\bigl((x_n+u_j)^{-s_j}-x_n^{-s_j}\bigr),
 \qquad\operatorname{Re}S>0.
\]
For $\operatorname{Re}S>m$ it is the inclusion--exclusion of $2^m$ separately
convergent spectra.

\begin{theorem}[Critical connected pole]\label{connected:germ}
There is a holomorphic function $\mathcal H_m$ on
$\operatorname{Re}S>-1$ such that
\begin{equation}\label{connected:pole}
 \mathcal C_m(\boldsymbol{s})=
 \left(\prod_js_j\right)
 \left(\frac{(-1)^mU}{S}+\mathcal H_m(\boldsymbol{s})\right).
\end{equation}
Its exact normalization is
\begin{align*}
 \mathcal H_m(\boldsymbol{s})={}&(-1)^mU
       \left(\zeta(1+S,a)-\frac1S\right)\\
 &+\sum_{n\ge0}\left[x_n^{m-1}\prod_jD_{s_j}(x_n,u_j)
                      -(-1)^mU x_n^{-1-S}\right].
\end{align*}
If $m\ge2$ and $U\ne0$, the unnormalized connected function has no jointly
holomorphic germ at the origin, although every non-tangent spectral line
has a regular restriction. On the equal line its derivatives below order
$m-1$ vanish, and
\[
 \left.\frac{d^{m-1}}{dt^{m-1}}\mathcal C_m(t,\ldots,t)\right|_{t=0}
 =\frac{(-1)^m(m-1)!}{m}U.
\]
\end{theorem}
\begin{proof}
Uniformly on compact parameter sets,
$D_s(x,u)=-u x^{-1-s}+O(x^{-2-\operatorname{Re}s})$.
The subtracted bracket is therefore $O(x^{-2-\operatorname{Re}S})$,
with the same summability after fixed derivatives introduce logarithmic
factors. This proves normal convergence. Restoring the subtracted leading
term gives the Hurwitz function, whose residue at $S=0$ is one.
The polynomial $\prod_js_j$ is not divisible by $S$ for $m\ge2$.
On a non-tangent line the quotient instead becomes a monomial of degree
$m-1$. The stated derivative follows by setting every $s_j=t$.
\end{proof}

For a multi-index $\boldsymbol\alpha\in\mathbb N_0^m$, put $N=|\boldsymbol\alpha|$
and
\[
 Q_r(x,u)=\frac{\log^{r+1}(x+u)-\log^{r+1}x}{r+1}.
\]
Use the convention
$\zeta(1+w,a)=1/w+\sum_{N\ge0}(-1)^N\gamma_N(a)w^N/N!$.

\begin{theorem}[All-index primitive and complete mass collision]
\label{connected:primitive}
The series
\begin{equation}\label{connected:F}
 \mathcal F_{\boldsymbol\alpha}(a;\boldsymbol u)
 =U\gamma_N(a)+\sum_{n\ge0}
 \left[x_n^{m-1}\prod_jQ_{\alpha_j}(x_n,u_j)
                  -\frac{U\log^N x_n}{x_n}\right]
\end{equation}
is normally and absolutely convergent. It is
$(-1)^{m+N}\partial_{\boldsymbol s}^{\boldsymbol\alpha}\mathcal H_m(0)$.
Its mixed mass derivative, taken before restriction to the diagonal, is
\begin{align}\label{connected:collision}
 &\left.\partial_{u_1}\cdots\partial_{u_m}
       \mathcal F_{\boldsymbol\alpha}(a;\boldsymbol u)
       \right|_{u_1=\cdots=u_m=u}\\
 &\qquad=\gamma_N(a+u)+(-1)^N
    \sum_{r=1}^{m-1}\binom{m-1}{r}(-u)^r
                           \zeta^{(N)}(r+1,a+u).\nonumber
\end{align}
Here $\zeta^{(N)}$ differentiates the order, not the shift.
\end{theorem}
\begin{proof}
Differentiation of $D_s$ gives
$\partial_s^rD_s|_0=(-1)^{r+1}Q_r$, and every subtracted bracket in
\eqref{connected:F} is $O(x^{-2}(1+\log^N x))$.
A mixed mass derivative replaces $Q_r$ by $\log^r(x+u)/(x+u)$.
At a complete collision the logarithms combine to $\log^N(x+u)$.
The finite identity
\[
 \frac{x^{m-1}}{(x+u)^m}
 =\sum_{r=0}^{m-1}\binom{m-1}{r}\frac{(-u)^r}{(x+u)^{r+1}}
\]
reduces every $r\ge1$ sum to a Hurwitz order derivative; the $r=0$
difference is $\gamma_N(a+u)-\gamma_N(a)$. This proves the result.
\end{proof}

Write $\mathcal F_m=\mathcal F_{(0,\ldots,0)}$. Here $\log\Gamma$ is the holomorphic logarithm in the right half-plane, real on the positive axis. At depth one,
$\mathcal F_1(a;u)=\log\Gamma(a)-\log\Gamma(a+u)$.
For positive real $u$, let
\[
 P_m(L)=\sum_{j=0}^{\lfloor(m-1)/2\rfloor}
          (-1)^j\binom m{2j+1}\pi^{2j}L^{m-2j-1}.
\]
Then the ordinary integral formula is
\begin{equation}\label{connected:gamma}
 \mathcal F_m(a;u,\ldots,u)
 =-\int_0^u t^{m-1}\psi(a+t)
                   P_m\left(\log\frac{t}{u-t}\right)\,dt.
\end{equation}
Indeed the unmarked mixed derivative is the negative divided difference
$-[u_1,\ldots,u_m](t^{m-1}\psi(a+t))$. Integrating over the mass box and
using Cauchy's divided-difference formula gives a contour integral of the
product of $m$ logarithms. The lower-minus-upper jump at equal masses is
$2\pi iP_m(\log(t/(u-t)))$. Its endpoint growth is a finite logarithmic
power, so the contour shrinks to the displayed absolutely convergent
integral. In particular $P_3(L)=3L^2-\pi^2$; the constant term cannot be omitted.

The all-index Gauss law is
\begin{align}\label{connected:gauss}
 &\sum_{h=0}^{q-1}\mathcal F_{\boldsymbol\alpha}
    \left(\frac{a+h}{q};\frac{\boldsymbol u}{q}\right)\\
 &\quad=q^{1-m}\left\{
 \sum_{\boldsymbol\beta\le\boldsymbol\alpha}
   \binom{\boldsymbol\alpha}{\boldsymbol\beta}
   (-\log q)^{N-|\boldsymbol\beta|}
       \mathcal F_{\boldsymbol\beta}(a;\boldsymbol u)
 +\frac{(-1)^NU(\log q)^{N+1}}{N+1}\right\},\nonumber
\end{align}
where $q\ge1$ is an integer. It follows by splitting the defining normalized
spectral series into residue classes, then subtracting the specified pole.

\paragraph{Normalization and scope.}
Mass collision is not the same operation as setting $S=0$ in the unnormalized
spectral expression. The mixed mass derivative is not the total derivative
of a one-variable diagonal function. These results are for this explicit
connected linear-factor family, not a general restriction theorem for singular
Stieltjes distributions. The accompanying full article proves the all-index
polylogarithmic expansion and Gamma-normalized antiderivatives and records
classical inputs, source-review limits, and independent diagnostic checks.
